\section{Síntesis y ampliación}

\subsection{Conclusiones centrales}

\begin{enumerate}
\item Los LLM enfrentan rendimientos marginales decrecientes de los datos; la robótica enfrenta un desierto de datos reales.
\item Los datos para robótica deben entenderse como una pirámide: datos reales recopilados en el dispositivo, datos de simulación o sintéticos y datos de internet o de personas.
\item El núcleo del Data Engine es un ciclo continuo: detectar fallas, completar los datos, entrenar, evaluar y volver a desplegar.
\item El valor de los datos no se mide en GB, sino por su autenticidad, su escasez, su verificabilidad y su rendimiento marginal.
\item La Data Recipe es una barrera de entrada implícita de la industria de datos.
\item El panorama de la industria robótica es una red simbiótica de empresas de modelos grandes, de modelos del mundo, de VLA, de cuerpos robóticos y de datos de simulación.
\item El destino final de los datos no es que desaparezcan, sino el paso de los datos estáticos a los entornos dinámicos y al autoaprendizaje.
\end{enumerate}

\subsection{Preguntas abiertas}

\begin{itemize}
\item Dentro de la pirámide de datos para robótica, ¿pueden los datos de simulación superar de verdad la sim-to-real gap?
\item ¿Quién controlará el ciclo cerrado de datos de la robótica: las empresas de cuerpos robóticos, las de modelos grandes o las de datos de simulación?
\item ¿Puede la Data Recipe convertirse en una barrera central, como lo es la arquitectura del modelo?
\item Cuando se fortalezca el autoaprendizaje de la IA, ¿la data factory se convertirá en una environment factory?
\end{itemize}

\subsection{Lecturas adicionales}

\begin{itemize}
\item Tesla FSD data engine: para entender el ciclo cerrado de datos recopilados en el dispositivo.
\item Artículos sobre Behavior benchmark, VLA y world model: para entender la relación entre la evaluación de robots y los modelos del mundo.
\item Materiales sobre datos sintéticos, simulación, domain randomization y sim-to-real: para entender los retos centrales de los datos para robótica.
\item EP137 AI for Math: para entender cómo un entorno verificable se convierte en una fuente de datos de alta calidad.
\end{itemize}

\begin{importantbox}{Juicio final}
El futuro de los datos no son «más rastreadores web», sino «mejores entornos». Quien logre construir entornos verificables, escalables y capaces de generar retroalimentación de fallas tomará la iniciativa en la siguiente etapa del entrenamiento de la inteligencia.
\end{importantbox}

\end{document}
